%==============================================================================
% The Illusion of Local: A Four-Layer Forensic Measurement of Privacy
% Boundaries in Consumer Local-LLM Serving Stacks
%
% Target venue : USENIX Security (Symposium on Usable Privacy and Security /
%                USENIX Security '26)
% Document class: article + usenix.sty -- two-column, 10pt, USENIX camera-ready
% Build        : pdflatex main && pdflatex main  (references are INLINE via
%                thebibliography; no .bib file needed)
%
% Style        : USENIX two-column camera-ready. Voice/structure conventions
%                inherited from the authors' ESWA D-26-11044 benchmark
%                paper (custom commands, finding boxes, cmark/xmark
%                tables, pgfplots with nodes near coords, etc.).
%
% ----------------------------------------------------------------------------
% NOTES
%  * Architecture diagrams and data charts that should be VECTOR/REPRODUCIBLE
%    are drawn natively in TikZ/pgfplots (no external image needed) so the
%    numbers in the charts stay tied to the measured data.
%  * Citations identify the publications and technical documents used here.
%==============================================================================
\documentclass[journal,tifs]{IEEEtran}

% usenix.sty already loads: mathptmx, fontenc, inputenc, pslatex, microtype,
% cite, url, xcolor, hyperref. Do NOT re-load these.

\usepackage{amsmath,amssymb,amsfonts}
\usepackage{graphicx}
\usepackage{booktabs}
\usepackage{multirow}
\usepackage{array}
\usepackage{colortbl}
\usepackage{balance}
\usepackage{rotating}
\usepackage{xcolor}
\usepackage{enumitem} % Required for customizing list environments
\usepackage{wrapfig}
\usepackage{xurl}
\usepackage{tikz}
\usepackage{pgfplots}
\pgfplotsset{compat=1.18}
\usetikzlibrary{shapes.geometric,arrows.meta,positioning,fit,backgrounds,calc,shadows}
\usepackage{subcaption}
\usepackage{siunitx}
\usepackage{mathtools}
\usepackage{bm}
\usepackage[most]{tcolorbox}
\usepackage{pifont}
\emergencystretch=1em
\usepackage{lineno}
\usepackage{doi}

% ---- Custom commands ----------
\newcommand{\cmark}{\ding{51}}
\newcommand{\xmark}{\ding{55}}
\newcommand{\hmark}{\ensuremath{\circ}}
\newcommand{\gguf}{\textsc{Gguf}}
\newcommand{\lms}{\textsc{Lm~Studio}}
\newcommand{\kvc}{KV-cache}
\newcommand{\tikzmark}[1]{\tikz[remember picture,overlay]\node(#1){};}

\newcommand{\ik}[1]{{\em\color{red}{[#1 ---ikram]}}}
\newcommand{\zc}[1]{{\em\color{blue}{[#1 ---zodiac]}}}
\newcommand{\revised}[1]{{\color{blue}#1}}
\newcommand{\artifacturl}[1]{\texttt{#1}}
% ---- Colors ------------------
\definecolor{vulncol}{RGB}{214,39,40}
\definecolor{fixcol}{RGB}{44,160,44}
\definecolor{partialcol}{RGB}{255,127,14}
\definecolor{netcol}{RGB}{31,119,180}
\definecolor{tableheader}{RGB}{240,240,240}
\definecolor{findingbg}{RGB}{255,250,230}
\definecolor{findingframe}{RGB}{200,160,40}
\usepackage{tikz}
\usetikzlibrary{arrows.meta,positioning,fit,calc}

\hypersetup{
	colorlinks=true,
	linkcolor=blue!60!black,
	citecolor=green!50!black,
	urlcolor=purple!70!black,
	pdfstartview={FitH}
}

\newtcolorbox{findingbox}{
	colback=findingbg,
	colframe=findingframe,
	boxrule=0.6pt,
	arc=2pt,
	left=4pt, right=4pt, top=3pt, bottom=3pt,
	fonttitle=\bfseries\small
} % <--- تم إضافة قوس الإغلاق هنا




\begin{document}

	%\title{LLAnalyzer: The Illusion of ``Local'': A Four-Layer Forensic Measurement of Privacy Boundaries in Consumer Local-LLM Serving Stacks}
    \title{The Illusion of Local Privacy:
Confidentiality Boundary Failures in Consumer LLM Serving Systems}

	% \author{
	% 	{\rm Youssef~Hamdi~Zafaan~Ibrahim} \\
	% 	Independent Researcher \\
	% 	{\tt youssefhamdi329@gmail.com}
	% 	\and
	% 	{\rm Mohammed~Khalaf~Salama} \\
	% 	Independent Researcher \\
	% 	{\tt mkhalafsalama@gmail.com}
 %        \and
	% 	{\rm Muhammad~Ikram} \\
	% 	Macquarie University \\
	% 	{\tt muhammad.ikram@mq.edu.au}
	% }
\author{}
% \vspace{-2cm}
	\maketitle
% \textcolor{red}{Usenix Security review is double blind, so comment the authors list.}
% \vsp
	\begin{abstract}

Local inference is often treated as a privacy guarantee because prompts need not leave the user's device. Local execution, however, specifies where computation occurs, not how serving software controls prompt data before, during, or after inference. We present \textbf{LLAnalyzer}, a measurement framework that traces prompt exposure across four confidentiality boundaries: model-artifact integrity, runtime lifetime, wrapper-managed persistence, and shared-serving isolation. We evaluate four open-weight model families and two consumer deployment platforms using structured artifact tests, memory forensics, persistent-artifact acquisition, and shared-serving measurements. Retained GGUF reports contain assertion-abort observations but no recorded SIGSEGV; no independent artifact supports the historical coverage-guided campaign total. Under the documented memory-acquisition path, post-inference prompt bytes remained recoverable from writable process mappings, including 11 of 12 sequential Nemotron tenants in each of two single-run observations. In the recorded LM~Studio configuration, sensitive-data logging produced 10/10 canary recoveries and disabling it produced none in the examined artifacts; six hook summaries show external connections but no canary in visible outbound payloads. In one unauthenticated shared-server process, slot save/restore produced 200/200 token-count equality checks. A separate retained continuation workload produced 126/200 KV-channel hits (Gemma 50/50, Qwen 50/50, Nemotron 25/50, Phi-4 1/50) and 200/200 memory-channel hits; its records are dated July 3--4, 2026. These observations show that local execution is a deployment property, not a complete prompt-confidentiality guarantee.

\end{abstract}

\begin{IEEEkeywords}
Large language models, information privacy, memory forensics, access control, side-channel analysis, security measurement.
\end{IEEEkeywords}
%================== Paper structure ===========    
% 1. Introduction (1.5 pages)

% 2. Confidentiality Boundary Framework
%    (Background + Threat Model) (1.5 pages)

% 3. LLAnalyzer
%    (Methodology) (2 pages)

% 4. Evaluation
%       Finding 1
%       Finding 2
%       Finding 3
%       Finding 4
%    (5 pages)

% 5. Discussion
%    (1 page)

% 6. Related Work
%    (1 page)

% 7. Conclusion
%    (0.5 page)
%================== Paper structure ===========
%==============================================================================
%==============================================================================
\section{Introduction}
\label{sec:introduction}
%==============================================================================

Large language models (LLMs) can now run entirely on personal computers through tools such as \emph{LM Studio}~\cite{lmstudio}, \emph{Ollama}~\cite{ollama}, and \texttt{llama.cpp}~\cite{llamacpp}. For privacy-sensitive workloads, the appeal is straightforward: prompts can be processed without sending them to a cloud inference provider. But what does ``local'' actually guarantee about a prompt after it reaches the serving software?

We find that this distinction matters. Keeping computation on-device determines where inference takes place; it says much less about what happens to a prompt while inference is being prepared, executed, and torn down. A single request may pass through model-loading code, HTTP and JSON parsers, chat-template processing, runtime memory, application storage, and a serving interface shared by multiple clients. These components make independent decisions about copying, retaining, persisting, and exposing prompt data. We therefore ask whether the privacy intuition associated with local inference survives this broader software lifecycle.

Research on LLM privacy has examined different parts of the problem: memorisation of training examples~\cite{carlini2023memorization,carlini2025}, forensic recovery of local prompts and conversation artifacts~\cite{murtuza2026,localnotprivate2026,xu2026mobilellmforensics}, and timing leakage or reconstruction from shared KV-cache state~\cite{wu2025iknowwhatyouasked,safekv2025,prefixwall2026,shadowcache2026,optileak2026}. These strands motivate a systems-level question: what happens to a user's prompt across the ordinary software stack of a consumer local-LLM deployment?

This leads us to a systems question: \emph{when inference stays on the user's machine, which components of the serving stack actually determine whether the prompt stays confidential?} Answering it requires more than searching for individual implementation bugs. If we treat the entire local stack as one trusted component, a failure observed at the serving interface is indistinguishable, conceptually, from one caused by runtime-memory retention or application-level persistence. It also becomes difficult to determine where a mitigation should be applied.

We instead study prompt confidentiality as a composition of explicit software boundaries. We develop \textbf{LLAnalyzer}, a measurement framework that separates a consumer local-LLM serving stack into four boundaries that we can exercise independently: \emph{(i)} an \textbf{Integrity Boundary}, covering model admission and loading; \emph{(ii)} a \textbf{Lifetime Boundary}, covering the lifetime of prompt data in runtime memory; \emph{(iii)} a \textbf{Persistence Boundary}, covering prompt material retained by wrapper applications; and \emph{(iv)} an \textbf{Isolation Boundary}, covering separation between clients of the serving interface. For each boundary, we inject controlled inputs, acquire boundary-specific evidence, and trace an observed failure back to the software layer in which it occurs.

Applying LLAnalyzer to four open-weight model families and two consumer deployment platforms produces a markedly uneven security picture. At the integrity boundary, retained structured reports show assertion aborts in malformed-input testing and no recorded SIGSEGV; the available evidence does not establish the historical AFL++ campaign total. The runtime boundary behaves differently. We recover prompt bytes after inference and trace the residue to multiple representations created during ordinary request processing and subsequently retained in allocator-managed memory. The two retained Nemotron twelve-tenant configurations are single-run observations and show the same 11/12 tenant-level outcome, so they do not identify a causal mitigation effect.

The failures are not confined to volatile memory. We find that consumer wrappers can extend prompt lifetime through plaintext persistence. At the serving interface, one unauthenticated loopback server and one Python process performed 200/200 token-count equality checks when a saved slot file was restored into another slot. The script did not use a second client, distinct API keys, an ownership control, or a wrong-key control; the result therefore demonstrates same-process state transfer under the tested configuration, not a cross-principal authorization bypass. A separate July-dated continuation workload observed 126/200 KV-channel hits and 200/200 memory-channel hits after restoration. We report these as boundary-specific observations and deployment-hardening guidance, without attributing them to a vendor decision. Timing vectors and the original analysis script are not retained, so timing separation is not used as a validated quantitative result.

These results change how we think about the privacy claim attached to local inference. Moving computation from a cloud provider onto a user's machine removes an important external trust relationship, but it does not establish what happens to prompt data inside the local serving stack. In the systems we study, confidentiality depends on the composition of model admission, runtime data lifetime, persistent storage, and client isolation. Local execution is therefore one condition in this trust model, rather than its endpoint.

\textbf{Contributions.}
Our work makes the following contributions:

\begin{itemize}

\item \textbf{A boundary-oriented model of local-LLM prompt
confidentiality.}
We formulate prompt confidentiality as a systems property spanning model admission, runtime lifetime, application persistence, and client isolation. This decomposition gives us explicit security objectives against which each layer can be measured.

\item \textbf{LLAnalyzer.}
We design and implement LLAnalyzer, a measurement framework for exercising these boundaries independently and attributing observed prompt-confidentiality failures to the software components responsible for them.

\item \textbf{An empirical study of confidentiality across the local
serving stack.}
Across four model families and two consumer deployment platforms, we find sharply different behaviour across boundaries: assertion-abort observations in retained GGUF reports, persistent post-inference prompt residue, configuration-dependent wrapper-level plaintext persistence, and same-process saved-state transfer under a shared-server configuration. We additionally report the limitations of the retained timing record rather than treating it as an independently validated result.

\item \textbf{Mechanisms, mitigations, and evidence limits.}
We trace the mechanisms underlying the observed runtime residue, compare the recorded sanitisation configurations descriptively, and report provenance and disclosure status conservatively. The retained records contain no vendor decision for the LM~Studio persistence or \texttt{llama-server} observations. The dated \texttt{DISCLOSURE\_PLAN.md} entry is a researcher-side scope determination. The executed June binaries were not hashed and are no longer available; later CPU-only rebuilds are retained for reference but are not claimed equivalent. Missing timing vectors, aggregate-only allocator records, and chronology limits are documented in the Supplement.
\end{itemize}

\noindent\textbf{Paper Structure.}
Section~\ref{sec:framework} defines our confidentiality boundaries and threat model. Section~\ref{sec:methodology} describes LLAnalyzer and our measurement methodology. Section~\ref{sec:evaluation} evaluates each boundary and investigates the mechanisms behind the observed failures. In Section~\ref{sec:related}, we position our findings against prior work, and conclude our work in Section~\ref{sec:conclusion}. 

\begin{figure*}[ht]
	\centering
	\includegraphics[width=0.85\textwidth]{figures/fig_llanalyzer_arch.pdf}
	\caption{Architecture and evidence map for LLAnalyzer. The concrete
		request path is shown above the four confidentiality boundaries:
		model-file integrity (B1), runtime lifetime (B2), wrapper-managed
		persistence (B3), and shared serving-state isolation (B4).}
	\label{fig:pb-llanalyzer}
\end{figure*} 

%==============================================================================
\section{Confidentiality Boundary Framework}
\label{sec:framework}
%==============================================================================

We begin by defining what prompt confidentiality means in a consumer
local-LLM deployment. We do not treat the deployment as a single trusted
environment. Instead, we follow a prompt through the serving stack and
identify the points at which different software components assume
responsibility for protecting it. This gives us four boundaries with
distinct security objectives and, importantly, distinct ways of failing.

%==============================================================================
\subsection{Prompt Lifecycle in Consumer LLM Serving Systems}
%==============================================================================

A typical local-LLM request begins in a consumer application such as
\emph{LM Studio} or \emph{Ollama} and eventually reaches an inference
engine such as \texttt{llama.cpp}. Between receiving the user's text and
returning a response, the stack loads the model, parses and transforms
the request, constructs the prompt, tokenizes it, manages inference
state, and may record conversation state or expose it through a serving
API.

Although this appears to the user as one local operation, the prompt does
not remain in one representation or under the control of one component.
It can be copied during parsing, transformed by a chat template, retained
in runtime memory, written to application storage, or associated with
state exposed through a serving interface. A confidentiality guarantee
at one of these stages therefore says little about the others.

Figure~\ref{fig:pb-llanalyzer} follows this lifecycle and shows where
LLAnalyzer places the corresponding confidentiality boundaries.

%==============================================================================
\subsection{Confidentiality Boundaries}
%==============================================================================

We define four boundaries according to the security decision being made
about prompt data.

\noindent\textbf{Integrity Boundary.}
Before inference begins, the model-loading pipeline decides whether an
input model artifact can safely enter execution. GGUF parsing, metadata
validation, tensor validation, and model initialisation form this
boundary. Its objective is to prevent malformed model artifacts from
compromising trusted model loading or subsequent execution.

\noindent\textbf{Lifetime Boundary.}
Once a request is processed, prompt data can exist simultaneously in
several runtime representations produced by parsing, template
construction, tokenisation, request processing, and memory allocation.
The lifetime boundary concerns whether those representations remain
recoverable after they are no longer required for inference.

\noindent\textbf{Persistence Boundary.}
Outside process memory, wrapper applications may retain conversation
histories, logs, caches, temporary files, and other application
artifacts. This boundary concerns whether prompt data are written to
persistent storage, for how long, and under what configuration.

\noindent\textbf{Isolation Boundary.}
When multiple clients interact with the same serving process, the
serving interface must keep their state separate. Authentication,
authorization, conversation-state management, cache sharing, and session
handling form this boundary. Its objective is to prevent one client from
learning prompt or conversation information belonging to another.

These boundaries are deliberately separated because satisfying one does
not imply satisfying another. A model may load safely while prompts
remain in memory; runtime memory may be sanitised while a wrapper writes
the same prompt to disk; and a locally protected process may still expose
another tenant's state through its API. We use this decomposition to
identify which guarantee is being tested in each experiment.

%==============================================================================
\subsection{Threat Model}
\label{sec:threat}
%==============================================================================

We consider adversaries that can exercise one of these boundaries without
already controlling the operating system or inference process. We assume
that the underlying operating system, hardware, and cryptographic
primitives operate correctly unless an experiment explicitly states
otherwise. Our goal is therefore not to model a fully compromised host,
but to determine what prompt information becomes available through the
ordinary privileges and interfaces surrounding a local-LLM deployment.

\noindent\textbf{${A}_1$ (Co-resident Local Adversary).}
${A}_1$ is an unrelated, unprivileged same-UID process with no assumed
root access, $\texttt{CAP\_SYS\_PTRACE}$, global Yama relaxation,
target-side $\texttt{PR\_SET\_PTRACER}$ permission, or parent/supervisor
relationship. Same UID therefore does not imply arbitrary access to
$\texttt{/proc/<pid>/mem}$; unauthorized access claims are made only when
directly demonstrated.

\noindent\textbf{$F$ (Authorized Forensic Collector).}
$F$ is the researcher-controlled collector used for memory measurement. It
uses a child or explicit-permission path and the documented
$\texttt{/proc/<pid>/maps}$ and $\texttt{/proc/<pid>/mem}$ reads. Its results
establish residual-data existence and conditional forensic recoverability,
not ${A}_1$ accessibility.

\noindent\textbf{${A}_2$ (Network Client).}
${A}_2$ is an authenticated client restricted to the published serving
interface. It has no direct access to process memory, the host filesystem,
or the underlying operating system. We use this adversary to test whether
the serving protocol preserves isolation between clients. The retained slot
and continuation workloads did not instantiate two such clients or distinct
credentials; they therefore do not provide an $A_2$ authorization result.

\noindent\textbf{${A}_3$ (Malicious Model Provider).}
${A}_3$ can supply a crafted or malformed GGUF artifact to an otherwise
unmodified inference engine. We use this capability to test whether the
model-loading pipeline rejects malformed artifacts without compromising
trusted execution.

Table~S4 in Supplementary Sec.~S1 connects each boundary
to the actor capable of challenging it and to the property we subsequently
measure. For Lifetime, $F$ is the measurement actor, whereas ${A}_1$ remains
the unauthorized co-resident adversary. Throughout, we distinguish (i)
residual-data existence, (ii) conditional forensic recoverability under
documented permissions, (iii) adversarial accessibility, and (iv) the
resulting confidentiality consequence.

%==============================================================================
\section{LLAnalyzer}
\label{sec:methodology}
%==============================================================================

LLAnalyzer operationalizes the four boundaries above as controlled
measurement experiments. Its central design principle is simple: when
testing one boundary, we vary the mechanism relevant to that boundary
while keeping the remainder of the serving configuration fixed. The
result is not a single vulnerability scanner, but a collection of
boundary-specific measurement harnesses connected by a common
experimental workflow.

%==============================================================================
\subsection{Design Goals}
%==============================================================================

We designed LLAnalyzer around three goals.

\noindent{\bf Boundary attribution.}
A recovered prompt is useful evidence only if we can determine how it
became exposed. LLAnalyzer therefore associates every measurement with a
specific boundary and collects evidence from the software component that
implements that boundary. This lets us distinguish, for example,
runtime-memory residue from wrapper persistence or protocol-level
shared-state exposure.

\noindent{\bf Methodological consistency.}
The boundaries cannot all be tested using the same instrumentation:
model admission requires parser testing, runtime lifetime requires memory
forensics, persistent storage requires artifact analysis, and tenant
isolation requires protocol-level experiments. We nevertheless impose
the same experimental sequence on each measurement so that the resulting
evidence answers the same question: did the target boundary satisfy its
security objective?

\noindent{\bf Reproducibility.}
We establish experimental ground truth using randomly generated UUID
canaries and controlled execution configurations. The retained
implementation, configurations, and supporting artifacts are described in
the Supplement. Source revision 721354f (full SHA
721354fbdfb7743e2be2183d918a3cdb9276c70f) is recorded in the project
reports; the binaries executed in June 2026 were not hashed and are no
longer available. An October 2026 CPU-only rebuild is retained for
reference, is not claimed byte-identical, and its equivalence to the
executed binary is unknown. LM~Studio runs used its bundled llama.cpp backend; the
GPU runs used a CUDA build whose revision is not recorded. Experiments
were run in June 2026 on the Table~S1 hardware, are reported from retained
records, and were not re-run for this revision.
For each experiment, LLAnalyzer generates
a fresh canary and embeds it in the submitted prompt, providing an
experiment-specific marker whose subsequent recovery can be attributed
to the originating request rather than inferred from surrounding content.

%==============================================================================
\subsection{Framework Architecture}
%==============================================================================

Figure~\ref{fig:pb-llanalyzer} shows how these measurement components map
onto the serving stack. For each boundary, LLAnalyzer follows five steps:
(\textit{i}) state the security objective;
(\textit{ii}) select the corresponding adversary;
(\textit{iii}) exercise the target boundary under a controlled
configuration;
(\textit{iv}) acquire evidence from that boundary; and
(\textit{v}) test the evidence against the security objective.

The instrumentation used in Step~(\textit{iv}) changes with the boundary;
the experimental logic does not. This gives us a common structure for
interpreting evidence obtained from otherwise different parts of the
serving system.

%==============================================================================
\subsection{Boundary-Oriented Measurement}
%==============================================================================Table~S5 in Supplementary
Sec.~S1 shows the
evidence source used for each boundary. During an experiment, we change
only the conditions needed to exercise the target boundary and preserve
the remaining configuration. Where a UUID canary is subsequently recovered,
LLAnalyzer uses the boundary-specific evidence to determine the software
path through which that disclosure occurred.
%%%%%%%%%%%%%%%%%%%%%%%%%%%%%%%%%%%%%%%%%%%%%%%%%%%%%%%%%%%%%%%%%%%%%%%%%%%%%%
%%%%%%%%%%%%%%%%%%%%%%%%%%%%%%%%%%%%%%%%%%%%%%%%%%%%%%%%%%%%%%%%%%%%%%%%%%%%%%
\subsection{Measurement Pipeline}
%%%%%%%%%%%%%%%%%%%%%%%%%%%%%%%%%%%%%%%%%%%%%%%%%%%%%%%%%%%%%%%%%%%%%%%%%%%%%%

Although each confidentiality boundary requires different
instrumentation, every LLAnalyzer experiment follows the same measurement
pipeline, shown in Figure~\ref{fig:pb-llanalyzer}.

\textit{First, establish ground truth.}
For each inference request $r$, LLAnalyzer generates a fresh UUIDv4 using
the operating system's cryptographically secure random-number generator
and embeds its canonical representation directly into the submitted
prompt as \texttt{CANARY:<uuid>}. Before inference begins, the harness
records
$\langle
r,\mathrm{UUID}_r,\mathrm{boundary},\mathrm{timestamp}
\rangle$
%
in the experiment manifest. This provides a deterministic mapping between
every injected marker and its originating request. A UUIDv4 contains 122 random bits, making an accidental exact match
negligibly probable. Before each experiment, LLAnalyzer additionally
checks that the newly generated canary is absent from the target evidence
source. An exact post-inference match can therefore be associated with
the corresponding injected request with negligible probability of
accidental collision.

\textit{Second, exercise the target boundary.}
LLAnalyzer executes the experiment under a controlled configuration,
isolating the boundary under evaluation while holding unrelated
experimental factors constant wherever possible. The specific condition
being exercised depends on the corresponding security objective.

\textit{Third, acquire boundary-specific evidence.}
The evidence source depends on the boundary: parser execution for the
Integrity Boundary, process memory for the Lifetime Boundary,
wrapper-managed artifacts for the Persistence Boundary, and
protocol-level observations for the Isolation Boundary.

\textit{Finally, test the boundary objective.}
LLAnalyzer searches the acquired evidence for exact UUID matches and maps
each recovered occurrence to its originating request through the
experiment manifest. A boundary failure is recorded when recovery
satisfies the predefined disclosure criterion for that boundary. The
precise success and failure criteria are specified with the corresponding
experiments in Section~\ref{sec:evaluation}.


%%%%%%%%%%%%%%%%%%%%%%%%%%%%%%%%%%%%%%%%%%%%%%%%%%%%%%%%%%%%%%%%%%%%%%%%%%%%%%
\subsection{Experimental Configuration}
%%%%%%%%%%%%%%%%%%%%%%%%%%%%%%%%%%%%%%%%%%%%%%%%%%%%%%%%%%%%%%%%%%%%%%%%%%%%%%

The experiments were run in June 2026 on the Table~S1 hardware and are
reported from retained records; they were not re-run for this revision.
Conditions differ by experiment. The headline CPU lifetime and slot
records use the script-declared CPU server path; the GPU records use a
CUDA build whose revision is not recorded. GPU scope is partial CUDA
offload only; no full-offload or non-\texttt{llama.cpp} serving stacks were
tested. Earlier WSL2 records are not used to broaden the headline claims.

Memory-acquisition discipline: the headline harness launches the target as a
child and reads writable \texttt{/proc/}\textless pid\textgreater\texttt{/mem}
regions listed in \texttt{/proc/}\textless pid\textgreater\texttt{/maps};
\linebreak
these opens use \texttt{PTRACE\_MODE\_ATTACH\_FSCREDS} without an explicit
\texttt{PTRACE\_ATTACH}. Retained records report Yama~1 with no root,
\texttt{CAP\_SYS\_PTRACE}, or global relaxation; GDB and GPU host scans are
separate from ${A}_1$ evidence.

We select four open-weight model families spanning different parameter
scales and architectural designs. All evaluated artifacts use GGUF
Q4\_K\_M quantisation and were obtained from the
\texttt{lmstudio-community} organization on Hugging Face. The retained
records do not provide a run-local model SHA-256 manifest, so this paper
does not claim byte-identity across every historical condition.
Table~S6 in Supplementary Sec.~S1 summarizes the evaluated models.




%%%%%%%%%%%%%%%%%%%%%%%%%%%%%%%%%%%%%%%%%%%%%%%%%%%%%%%%%%%%%%%%%%%%%%%%%%%%%%
\subsection{Experimental Controls and Statistical Analysis}
%%%%%%%%%%%%%%%%%%%%%%%%%%%%%%%%%%%%%%%%%%%%%%%%%%%%%%%%%%%%%%%%%%%%%%%%%%%%%%

We use controlled comparisons to separate boundary-specific effects from
changes elsewhere in the serving stack. Within a retained experiment family,
the hardware, operating system, model artifact, software revision, and
inference parameters are held fixed unless the corresponding variable is
under evaluation; this is the basis for the slot, continuation, injection,
KV-overhead, and D4 comparisons. Across families, however, the retained
records are heterogeneous: the twelve-tenant rows are Nemotron-only, the
multilayer diagnostics are Phi-4-only, and the GPU-v5 campaign is a
Nemotron CUDA configuration. These observations are therefore interpreted
descriptively rather than as a hardware- or binary-equivalent factorial
design. When comparing configurations, we use identical prompt inputs and
model artifacts wherever the retained protocol permits.

For each boundary, we define an explicit success and failure criterion
and apply it to the retained records. A disclosure is recorded only when
an exact UUID match satisfies the corresponding boundary-specific
criterion. We report descriptive proportions and denominators; Wilson
intervals are used only where the retained trial unit supports them.
Continuous measurements, including inference latency and recovered-memory
volume, are summarized descriptively. No paired significance test is
claimed because the retained records do not provide a complete matched
analysis for the relevant comparisons.


% %%%%%%%%%%%%%%%%%%%%%%%%%%%%%%%%%%%%%%%%%%%%%%%%%%%%%%%%%%%%%%%%%%%%%%%%%%%%%%
% \subsection{Reproducibility}
% %%%%%%%%%%%%%%%%%%%%%%%%%%%%%%%%%%%%%%%%%%%%%%%%%%%%%%%%%%%%%%%%%%%%%%%%%%%%%%

% LLAnalyzer records an experiment manifest for every run containing the
% hardware and software configuration, model identity and SHA-256 digest,
% inference parameters, generated canaries, timestamps, and
% boundary-specific measurement settings. We retain the corresponding
% forensic evidence and analysis outputs required to reconstruct each
% reported result.

% We release the LLAnalyzer implementation, experiment configurations,
% analysis scripts, patches, and supporting artifacts to support independent
% reproduction and extension of our measurements. The artifacts are
% available at
% \url{https://github.com/0xzodiac/Local-LLMs-Forensics}.

%==============================================================================
\section{Evaluation}
\label{sec:evaluation}
%==============================================================================

We use LLAnalyzer to evaluate prompt confidentiality across the four
boundaries defined in Section~\ref{sec:framework}. We organize the
evaluation around research questions rather than individual measurement
tools. For each question, we state the security property under test,
describe the experimental design, present the resulting evidence, and
interpret what the observations establish about the corresponding
boundary.

Our evaluation addresses four research questions:

\begin{itemize}[left=0pt]

\item \textbf{RQ1. Integrity Boundary.}
Can malicious model artifacts bypass model-admission controls or trigger
memory-safety failures before inference begins?

\item \textbf{RQ2. Lifetime Boundary.}
Does prompt plaintext remain recoverable from runtime memory after
inference, and if so, what mechanisms determine its persistence and
mitigation?

\item \textbf{RQ3. Persistence Boundary.}
Do consumer wrapper applications extend prompt lifetime by retaining
prompt data in persistent local artifacts?

\item \textbf{RQ4. Isolation Boundary.}
Can saved state and subsequent continuation be observed across slot
identifiers in one shared serving process? Cross-principal authorization,
distinct credentials, ownership controls, and remotely observable side
channels are explicit untested boundaries in the retained experiments.

\end{itemize}

Together, these questions test our central hypothesis: 
%
% \begin{quote}
\emph{Local execution alone does not establish prompt confidentiality;
confidentiality depends on the security properties enforced across the
software boundaries through which prompt data pass.} 
% \end{quote}
%
We evaluate the boundaries in prompt-lifecycle order. We begin with model
admission (RQ1), then follow prompt data through runtime memory (RQ2) and
wrapper-managed persistent storage (RQ3), before examining information
separation between clients of a shared serving instance (RQ4). For each
experiment, we isolate the target boundary while holding unrelated
experimental factors constant wherever possible, following the controls
described in Section~\ref{sec:methodology}.


%==============================================================================
\subsection{RQ1. Integrity Boundary}
\label{sec:rq1}
%==============================================================================

Before a prompt can be processed, the serving framework must parse and
initialize the supplied GGUF model artifact. A failure at this stage can
invalidate assumptions made by later confidentiality measurements: if a
malformed model compromises the serving process before inference begins,
subsequent prompt disclosures may be consequences of an already
compromised execution environment. We therefore ask: \emph{Can a maliciously crafted GGUF artifact bypass parser validation or trigger memory-safety failures before inference begins?}


{\bf Experimental Design.}
We exercise the model-admission path with retained structured malformed
GGUF artifacts targeting tensor metadata, architecture identifiers,
dimensions, allocation sizes, offsets, RoPE parameters, and header
fields. The retained report families use different corpora and overlap;
we therefore keep their denominators separate. A separate historical
coverage-guided campaign is mentioned in older manuscript text, but no
fuzzer log, queue, executable, coverage record, or campaign manifest is
retained. It is not used as a quantitative result here.


{\bf Evidence.}
Table~\ref{tab:model_boundary_results} summarizes the retained reports.
The cross-model report describes 96 structured tests with seven SIGABRT
assertion aborts, 36 clean loads, 53 rejections, and zero SIGSEGV. The
surgical and tensor reports document additional assertion-abort outcomes,
including one non-deterministic tensor case, and no recorded SIGSEGV. These
families overlap and do not form a single deduplicated crash rate. The
reports describe four architectures, but their fourth column is labelled
`phi3` while the retained test scripts use the `phi4` key; the raw outputs
needed to establish that identity are not retained. No retained artifact
establishes the historical AFL++ execution total.


\begin{table}[!t]
\centering
\caption{Integrity-boundary evaluation using retained structured malformed
GGUF reports. Report families overlap; counts are not pooled into one crash
rate.}
\label{tab:model_boundary_results} 

\scriptsize
\begin{tabular}{@{}p{0.31\columnwidth}p{0.31\columnwidth}p{0.25\columnwidth}@{}}
\toprule
\textbf{Metric} & \textbf{Observation} & \textbf{Outcome} \\
\midrule
Cross-model structured tests & 96 & Reported \\
Model families             & 4                 & Reported   \\
Additional report families & -- & Overlapping \\
\midrule
SIGABRT assertion aborts   & 7 in cross-model report & Observed \\
Additional SIGABRT reports & 9 plus 1 non-deterministic & Observed \\
SIGSEGV                    & 0 recorded          & Not observed \\
Malformed model loaded     & Not established     & No retained positive \\
Fuzzer execution total     & Not retained        & Unverifiable \\
\bottomrule
\end{tabular}
\end{table}



{\bf Analysis.}
The retained evidence supports a narrower statement: malformed-input
testing produced assertion-abort process terminations in some report
families and no recorded SIGSEGV. It does not support a global parser-crash
zero, a unified crash rate, or the historical campaign total. The later
runtime and persistence observations were collected from serving processes
that completed their respective setup, but this does not turn the GGUF
results into a universal parser-security claim.


{\bf Scope of the Integrity Result.}
The reports do not establish universal GGUF-parser robustness. Their
corpora, denominators, and outcomes differ, and generated variants and
fuzzer artifacts are not retained for deduplication. The supported result
is an observation about the retained tests: assertion aborts occurred,
SIGSEGV was not recorded, and the available evidence cannot establish the
unretained campaign total or a complete memory-safety classification.


{\bf Takeaway.}
The integrity evidence is mixed rather than cleanly negative: retained
structured reports contain SIGABRT assertion aborts, while no retained
report records SIGSEGV. We next examine confidentiality observations that
arise in runtime memory, wrapper artifacts, and shared serving state.

%%%%%%%%%%%%%%%%%%%%%%%%%%%%%%%%%%%%%%%%%%%%%%%%%%%%%%%%%%%%%%%%%%%%%%%%%%%%%%
\subsection{RQ2. Lifetime Boundary}
\label{sec:rq2}
%%%%%%%%%%%%%%%%%%%%%%%%%%%%%%%%%%%%%%%%%%%%%%%%%%%%%%%%%%%%%%%%%%%%%%%%%%%%%%

We next examine what happens to prompt data after the serving process has
completed its setup. During inference, a prompt passes through
multiple runtime representations, including formatted templates, JSON
request structures, token buffers, temporary application objects, and
allocator-managed memory. Although these representations are transient
from the application's perspective, their underlying contents may
outlive the request that created them.

This section therefore addresses the following research question: \emph{Does runtime memory preserve prompt confidentiality after
inference, or does prompt plaintext remain recoverable beyond execution
completion?}  We address this question through four progressively deeper
sub-questions. We first establish whether plaintext remains recoverable
after inference (RQ2.1), trace the mechanisms responsible for its
persistence (RQ2.2), evaluate whether runtime sanitisation can reduce
this exposure at acceptable cost (RQ2.3), and finally examine whether
GPU-assisted execution changes these confidentiality properties (RQ2.4).

%%%%%%%%%%%%%%%%%%%%%%%%%%%%%%%%%%%%%%%%%%%%%%%%%%%%%%%%%%%%%%%%%%%%%%%%%%%%%%
\subsubsection{RQ2.1: Does Prompt Plaintext Persist After Inference?}
\label{sec:rq2_plaintext}
%%%%%%%%%%%%%%%%%%%%%%%%%%%%%%%%%%%%%%%%%%%%%%%%%%%%%%%%%%%%%%%%%%%%%%%%%%%%%%

We examine runtime state after inference and ask whether prompt plaintext
remains recoverable from process memory. The acquisition path is a
researcher-controlled parent process reading writable mappings of its
serving child; it is not a claim about an arbitrary same-UID or remote
attacker.

%%%%%%%%%%%%%%%%%%%%%%%%%%%%%%%%%%%%%%%%%%%%%%%%%%%%%%%%%%%%%%%%%%%%%%%%%%%%%%

{\bf Evidence.} Using the UUID-based forensic methodology described in
Section~\ref{sec:methodology}, $F$ acquires a post-inference memory
snapshot. Each prompt contains a fresh UUID canary, so an exact match can
be attributed to its originating request. In two retained Nemotron
twelve-tenant single-run configurations, canaries from 11 of 12 tenants were recovered
from writable, non-GGUF mappings. The corresponding exact byte-occurrence
totals were 110 copies and 10\,834~KB in one record, and 78 copies and
3\,352~KB in the other. These are byte-occurrence and printable-cleartext
measurements, not counts of independent secrets or attacker sessions.

\begin{table}[!t]
    \centering
    \caption{Retained Nemotron twelve-tenant runtime-memory observations.
    Each row is a single run; copy counts are exact byte occurrences in the
    scanned mappings and printable-cleartext volume is a scan metric.}
    \label{tab:residue}
    \footnotesize
    \setlength{\tabcolsep}{4pt}
    \begin{tabular}{@{}lccc@{}}
        \toprule
        Configuration & Tenants recovered & Copies & Cleartext (KB) \\
        \midrule
        Unpatched & 11/12 & 110 & 10\,834 \\
        Sanitised & 11/12 & 78 & 3\,352 \\
        \bottomrule
    \end{tabular}
\end{table}



%%%%%%%%%%%%%%%%%%%%%%%%%%%%%%%%%%%%%%%%%%%%%%%%%%%%%%%%%%%%%%%%%%%%%%%%%%%%%%

{\bf Analysis.} The two experiments establish complementary properties of runtime
residue. First, post-inference memory inspection shows that prompt
plaintext survives response generation: every measured inference leaves
multiple recoverable copies in process memory. The small variation in
copy count does not change the underlying observation that plaintext
recovery occurs consistently after inference. Second, the sequential-tenant experiment shows that this residue is not
confined to the most recent request. Canaries associated with earlier
tenants remain recoverable after subsequent requests have completed.
Prompt lifetime therefore extends beyond both response generation and
the execution of later inference requests within the same serving
process. These observations expose a mismatch between application-level and
memory-level lifetime: plaintext remains conditionally recoverable under
the documented $F$ path.

%%%%%%%%%%%%%%%%%%%%%%%%%%%%%%%%%%%%%%%%%%%%%%%%%%%%%%%%%%%%%%%%%%%%%%%%%%%%%%

{\bf Validating Runtime Residue.} We perform several checks to establish that the recovered plaintext
originates from the serving runtime rather than from LLAnalyzer's
acquisition procedure. First, prompt instances are tracked using fresh UUID canaries, enabling
recovered strings to be attributed to specific inference requests
through exact matching. Second, recovery is reproducible across
independent executions: although the precise number of copies may vary
slightly, residual prompt plaintext is recovered consistently after
inference. Third, the retained four-model slot, continuation, injection,
KV-overhead, and D4 families all provide model-labelled observations,
although their logical protocols and measurement channels differ. The two
twelve-tenant rows themselves are Nemotron-only, and the multilayer-v2
diagnostics are Phi-4-only; neither should be read as a balanced four-model
runtime-memory baseline. Most importantly, the sequential-tenant experiment provides temporal
provenance for the recovered residue. Canaries associated with earlier
tenants remain recoverable after subsequent inference requests have
completed, demonstrating that the recovered plaintext represents
historical runtime state rather than merely the currently executing
request. Together, these checks support attribution of the recovered plaintext
to post-inference state retained within the serving process rather than
to content introduced by LLAnalyzer's acquisition procedure.

%%%%%%%%%%%%%%%%%%%%%%%%%%%%%%%%%%%%%%%%%%%%%%%%%%%%%%%%%%%%%%%%%%%%%%%%%%%%%%

{\bf Takeaway.}
Prompt lifetime extends beyond inference completion in the evaluated
serving environment. Multiple plaintext representations remain
recoverable after response generation, and historical prompt contents
survive subsequent inference requests within the same serving process.
Runtime confidentiality therefore cannot be inferred from
application-level request completion alone.

% Having established that prompt plaintext persists after inference, 
We next investigate where these residual copies originate and why their
contents remain recoverable.

%%%%%%%%%%%%%%%%%%%%%%%%%%%%%%%%%%%%%%%%%%%%%%%%%%%%%%%%%%%%%%%%%%%%%%%%%%%%%%
\subsubsection{RQ2.2: Why Does Prompt Residue Persist?}
\label{sec:rq2_mechanism}
%%%%%%%%%%%%%%%%%%%%%%%%%%%%%%%%%%%%%%%%%%%%%%%%%%%%%%%%%%%%%%%%%%%%%%%%%%%%%%

RQ2.1 established that prompt plaintext remains recoverable after
inference and can persist across successive tenant requests within the
same serving process. We next investigate the mechanisms responsible
for extending prompt lifetime beyond the visible execution of an
inference request. Unlike conventional memory-forensics studies that identify surviving
objects without necessarily explaining their origin, LLAnalyzer
reconstructs the provenance of UUID-tagged prompt copies through
debugger-assisted runtime inspection. Rather than treating each
recovered copy as an independent artefact, we trace how a single user
prompt propagates through successive software layers during normal
inference execution.

%%%%%%%%%%%%%%%%%%%%%%%%%%%%%%%%%%%%%%%%%%%%%%%%%%%%%%%%%%%%%%%%%%%%%%%%%%%%%%

{\bf Evidence.}
Debugger-assisted provenance analysis follows the prompt throughout the
inference pipeline, including HTTP request decoding, JSON parsing,
chat-template rendering, inference request construction, token
generation, and response teardown. Across the evaluated executions,
recovered UUID-tagged prompt data are associated with three recurring
mechanisms:
% \begin{enumerate}
\textit{(1). Dynamic \texttt{std::string} growth.}
During prompt construction and template expansion, \texttt{std::string}
objects may repeatedly increase their capacity. When reallocation
occurs, the previous character buffer is released back to the allocator
without necessarily being overwritten immediately, allowing its
previous plaintext contents to remain recoverable.
\textit{(2). Transient runtime objects.} 
Intermediate processing stages, including JSON parsing, chat-template
construction, request serialisation, and related request-processing
operations, generate short-lived runtime objects containing copies or
derived representations of the original prompt. Although these objects
are destroyed after their immediate use, their underlying memory can
continue to contain recoverable prompt data. 
\textit{(3). Allocator-managed memory reuse.}
After object destruction, released buffers enter allocator-managed
reuse mechanisms, including per-thread caches and arena-managed regions
~\cite{glibc_arena}. These mechanisms retain released memory for subsequent
allocation rather than immediately clearing its previous contents,
thereby extending the period during which prompt plaintext can remain
recoverable. 
%
% \end{enumerate}
%
Collectively, these mechanisms account for the provenance of the
multiple prompt representations observed in RQ2.1.

%%%%%%%%%%%%%%%%%%%%%%%%%%%%%%%%%%%%%%%%%%%%%%%%%%%%%%%%%%%%%%%%%%%%%%%%%%%%%%

{\bf Analysis.}
The provenance analysis shows that prompt persistence is not caused by a
single programming error or forgotten deallocation. Instead, it emerges
from the interaction of software components that create multiple
plaintext representations during ordinary inference.

Importantly, the analysis separates prompt \emph{duplication} from
\emph{retention}. Multiple independent representations are created
before memory enters glibc's reuse structures through JSON parsing,
chat-template expansion, request construction, serialisation, and
repeated \texttt{std::string} reallocations. After the corresponding
objects are released, allocator behaviour extends the period during
which their underlying contents remain recoverable. Application-level data handling therefore explains why multiple prompt
representations exist, whereas allocator-managed memory reuse explains
why released representations may survive beyond their intended
application lifetime. The distinction is important because these two
mechanisms imply different mitigation requirements.

If allocator retention alone were responsible for the observed residue,
clearing allocator-managed memory would be sufficient to restore
confidentiality. The debugger traces instead show that sensitive
plaintext is propagated through several stages of the inference
pipeline before those allocations are released. Effective sanitisation
must therefore account for the lifecycle of prompt-derived
representations across the serving stack rather than targeting only
allocator state.

%%%%%%%%%%%%%%%%%%%%%%%%%%%%%%%%%%%%%%%%%%%%%%%%%%%%%%%%%%%%%%%%%%%%%%%%%%%%%%

{\bf Root-Cause Attribution.}
The debugger traces identify two successive mechanisms:
application-level prompt duplication followed by allocator-level
retention. Application processing creates the plaintext
representations; allocator behaviour subsequently extends their
recoverable lifetime. We therefore attribute prompt proliferation to
application-level data handling and persistence to subsequent memory
retention.

%%%%%%%%%%%%%%%%%%%%%%%%%%%%%%%%%%%%%%%%%%%%%%%%%%%%%%%%%%%%%%%%%%%%%%%%%%%%%%

{\bf Takeaway.}
Prompt residue results from the combined effect of application-level
duplication and subsequent memory retention. Prompt data are copied
across multiple stages of request processing, while allocator behaviour
allows some of those representations to remain recoverable after their
application objects have been released. Runtime prompt confidentiality
is therefore a lifecycle problem spanning multiple software components,
rather than an isolated property of the memory allocator.

% Having established how prompt plaintext is duplicated and retained, 
\textit{Next}, we investigate whether practical runtime mitigations can interrupt
this propagation chain and restore prompt confidentiality after
inference.

%%%%%%%%%%%%%%%%%%%%%%%%%%%%%%%%%%%%%%%%%%%%%%%%%%%%%%%%%%%%%%%%%%%%%%%%%%%%%%

{\bf Cross-model residue (D4).}
As a separate runtime-residue workload, we switched models within one
serving process and inspected the writable mappings before and after the
swap. The retained family contains 200 timestamped primary records: 50
Gemma-to-Phi-4, 50 Qwen-to-Gemma, 50 Nemotron-to-Qwen, and 50
Phi-4-to-Nemotron trials. The process identity is unchanged in all 200
records, and the canary from model~A remains present after the swap in all
200 records. A separately named debug JSON is not counted as an additional
trial. These are byte-occurrence observations in scanned mappings; the
records do not establish a model-generated disclosure or a cross-principal
access path.


%%%%%%%%%%%%%%%%%%%%%%%%%%%%%%%%%%%%%%%%%%%%%%%%%%%%%%%%%%%%%%%%%%%%%%%%%%%%%%
\subsubsection{RQ2.3: Can Runtime Sanitisation Reduce Residual
Recoverability at Acceptable Cost?}
\label{sec:rq2_mitigation}
%%%%%%%%%%%%%%%%%%%%%%%%%%%%%%%%%%%%%%%%%%%%%%%%%%%%%%%%%%%%%%%%%%%%%%%%%%%%%%

%RQ2.1 established that prompt plaintext remains recoverable after
%inference, while RQ2.2 showed that this persistence arises from multiple
%runtime representations created during normal request processing and
%subsequently retained in memory. 
We ask whether runtime sanitisation reduces residual prompt data and
conditional forensic recoverability at acceptable inference cost.

%%%%%%%%%%%%%%%%%%%%%%%%%%%%%%%%%%%%%%%%%%%%%%%%%%%%%%%%%%%%%%%%%%%%%%%%%%%%%%

{\bf Evidence.}
The two retained Nemotron twelve-tenant rows are single-run observations collected
under different executable paths and runtime configurations. Both recover
canaries from 11 of 12 tenants, with 110 versus 78 exact byte-occurrence
copies and 10\,834 versus 3\,352~KB of reported printable cleartext. The
rows therefore describe residual exposure under two configurations; they
do not identify a causal sanitisation effect or a confidence interval.
We separately measure KV-cache zeroisation overhead using 50 requests per
model. Table~\ref{tab:overhead} reports the retained mean latencies and
their arithmetic differences. These are descriptive workload observations,
not equivalence tests or universal performance bounds.

\begin{table}[!t]
    \centering
    \caption{Single-run Nemotron runtime-memory observations after twelve
    sequential tenants. Copy counts are exact byte occurrences in scanned
    mappings and printable-cleartext volume is a recovered-memory metric, not
    an exact count of secret-bearing bytes. Both rows recover canaries from
    11/12 tenants; the protocols are not a causal matched comparison.}
    \label{tab:blast}
    \footnotesize
    \begin{tabular*}{\columnwidth}{@{}l@{\extracolsep{\fill}}lccr@{}}
        \toprule
        Configuration & Tenants & Copies & Printable cleartext (KB) \\
        \midrule
        Unpatched & 11/12 & 110 & 10\,834 \\
        Sanitised & 11/12 & \textbf{78} & \textbf{3\,352} \\
        \bottomrule
    \end{tabular*}
\end{table}



\begin{table}[!t]
    \centering
    \caption{Descriptive KV-cache zeroisation overhead ($n=50$ requests per
    model). Mean request latency is reported in seconds; no equivalence
    margin or formal equivalence test was specified.}
    \label{tab:overhead}
    \footnotesize
    \begin{tabular}{@{}lrrr@{}}
        \toprule
        Model & Patched & Unpatched & $\Delta$ (s) \\
        \midrule
        Nemotron-3-Nano-4B      & 2.822 & 2.838 & $-0.016$ \\
        Qwen3.5-9B              & 4.705 & 4.692 & $+0.013$ \\
        Gemma-4-E4B-it          & 2.908 & 2.891 & $+0.017$ \\
        Phi-4-reasoning-plus    & 8.422 & 8.403 & $+0.020$ \\
        \bottomrule
    \end{tabular}
\end{table}

%%%%%%%%%%%%%%%%%%%%%%%%%%%%%%%%%%%%%%%%%%%%%%%%%%%%%%%%%%%%%%%%%%%%%%%%%%%%%%

{\bf Analysis.}
The observations distinguish residual magnitude from tenant-level
recoverability. The second row has fewer observed byte occurrences and a
smaller printable-cleartext metric, but both rows recover canaries from
11/12 tenants. Because the protocols differ in binary path, thread and
offload settings, startup and cleanup, and patch state, the retained rows
cannot support a causal percentage reduction. They instead motivate
matched repeated experiments as future work.

The latency arrays provide a separate descriptive result. The retained
mean differences are $-0.016$~s (Nemotron), $+0.013$~s (Qwen3.5),
$+0.017$~s (Gemma), and $+0.020$~s (Phi-4), corresponding to
$-0.5515$\%, $+0.2684$\%, $+0.5999$\%, and $+0.2324$\% when computed
from the retained means. No inferential or equivalence conclusion is made.

%%%%%%%%%%%%%%%%%%%%%%%%%%%%%%%%%%%%%%%%%%%%%%%%%%%%%%%%%%%%%%%%%%%%%%%%%%%%%%

{\bf Residual Exposure Analysis.}
The remaining exposure provides evidence about where sanitisation must
operate. RQ2.2 showed that prompt-derived plaintext can appear during
template expansion, JSON processing, request serialisation,
\texttt{std::string} growth, and subsequent allocator-managed retention.
A mitigation that clears only selected structures therefore protects
only the copies passing through those sanitised paths. This distinction also affects how mitigation effectiveness should be
measured. Copy count and recovered byte volume quantify the
\emph{magnitude} of residual exposure, whereas the number of tenants
whose prompts remain recoverable captures whether historical
confidentiality has been restored. Under the latter criterion, the two
retained configurations leave the same 11/12 tenant outcome.
Restoring the evaluated lifetime boundary therefore requires
lifecycle-wide handling of prompt-bearing representations, such that no
recoverable copy remains after its intended use. The results do not show
that runtime sanitisation is inherently incapable of achieving this
goal; rather, they show that the evaluated sanitisation coverage does
not yet encompass every surviving prompt representation.

%%%%%%%%%%%%%%%%%%%%%%%%%%%%%%%%%%%%%%%%%%%%%%%%%%%%%%%%%%%%%%%%%%%%%%%%%%%%%%

{\bf Takeaway.}
The retained twelve-tenant records show differences in raw copy and
printable-cleartext metrics, but both are single-run observations with
11/12 tenant recovery and confounded protocols. They support a descriptive
warning that selected sanitisation paths did not eliminate historical
recovery under the tested configurations, not a causal effect size. The KV
overhead arrays likewise support model-specific descriptive latency
differences only.

We next examine whether these runtime confidentiality properties change
when model execution is accelerated through GPU offloading.
%%%%%%%%%%%%%%%%%%%%%%%%%%%%%%%%%%%%%%%%%%%%%%%%%%%%%%%%%%%%%%%%%%%%%%%%%%%%%%



%%%%%%%%%%%%%%%%%%%%%%%%%%%%%%%%%%%%%%%%%%%%%%%%%%%%%%%%%%%%%%%%%%%%%%%%%%%%%%
\subsubsection{RQ2.4: Does GPU Execution Fundamentally Change Runtime
Confidentiality?}
\label{sec:rq2_gpu}
%%%%%%%%%%%%%%%%%%%%%%%%%%%%%%%%%%%%%%%%%%%%%%%%%%%%%%%%%%%%%%%%%%%%%%%%%%%%%%

We now examine whether GPU-assisted inference changes the host-memory
residue observed in the lifetime experiments.

A plausible hypothesis is that moving model computation to the GPU also
moves sensitive prompt state away from host memory. The test below is
limited to the recorded partial-offload configuration and host-memory
scan.

%%%%%%%%%%%%%%%%%%%%%%%%%%%%%%%%%%%%%%%%%%%%%%%%%%%%%%%%%%%%%%%%%%%%%%%%%%%%%%

{\bf Evidence.}
The retained v5 records contain 50 timestamped primary trials of a
Nemotron CUDA configuration. Immediately after one request, the recorded
	\texttt{canary\_hit\_count} is 9 in 47 trials and 10 in 3 trials. After two further
	requests, the \texttt{extra\_tenants[*].hits} value is 16 in every primary record.
	After process termination, the retained \texttt{vram\_after\_shutdown} field is
the tuple $(101,8192)$ in those 50 records. These fields are mapping/copy
measurements, not counts of successful attacks; the raw records do not
retain the previously described idle row, and they do not establish direct
device-memory recovery.

\begin{table}[!t]
    \centering
    \caption{\small Host-memory observations in the retained 50-trial CUDA
    record. The fields are measurement-specific and are not a binary
    exploit-success oracle.}
    \label{tab:vram}
    \scriptsize
    \begin{tabular}{@{}p{0.29\columnwidth}p{0.25\columnwidth}p{0.34\columnwidth}@{}}
        \toprule
        Observation point & Retained field & Interpretation \\
        \midrule
        After one request
            & 9 (47/50), 10 (3/50)
            & \texttt{canary\_hit\_count} distribution \\
        After two additional requests
            & 16 in all 50
            & extra-tenant hits field \\
        After server termination
            & $(101,8192)$ in all 50
            & \texttt{vram\_after\_shutdown} mapping tuple \\
        \bottomrule
    \end{tabular}
    \vspace{-0.5cm}
\end{table}

%%%%%%%%%%%%%%%%%%%%%%%%%%%%%%%%%%%%%%%%%%%%%%%%%%%%%%%%%%%%%%%%%%%%%%%%%%%%%%

{\bf Analysis.}
The retained record shows host-memory observations under partial CUDA
offload, including a 9/10-hit distribution after one request and a value
of 16 in the extra-tenant field after two further requests. The shutdown
tuple is a mapping observation, not proof that all memory recovery was
zero. Prompt construction and request processing remain host-side in this
configuration, but the records do not establish direct device-memory
contents or physical-page erasure.

%%%%%%%%%%%%%%%%%%%%%%%%%%%%%%%%%%%%%%%%%%%%%%%%%%%%%%%%%%%%%%%%%%%%%%%%%%%%%%

{\bf Host-Side Attribution.}
The provenance results from RQ2.2 help explain why GPU offloading does
not remove the observed residue. The recovered prompt representations
originate from host-side stages such as JSON parsing, template expansion,
request construction, tokenisation, \texttt{std::string} operations, and
subsequent allocator-managed retention. These operations occur before or
around GPU execution and therefore remain relevant even when substantial
model computation is offloaded. This attribution also defines the scope of
our result. The retained 50-trial GPU-v5 campaign directly evaluates
Nemotron-3-Nano-4B under partial CUDA offloading; it does not establish that
every fully offloaded model or alternative GPU-serving architecture exhibits
identical persistence. A historical report separately describes a Qwen3.5
GPU reproduction, but no linked primary record places that observation in
the retained GPU-v5 denominator. Nevertheless, within the
evaluated serving architecture, relocating model computation does not
remove the host-side software paths responsible for the recovered
plaintext. GPU offloading therefore cannot substitute for explicit
lifecycle management of prompt-bearing host memory in this configuration.

%%%%%%%%%%%%%%%%%%%%%%%%%%%%%%%%%%%%%%%%%%%%%%%%%%%%%%%%%%%%%%%%%%%%%%%%%%%%%%

{\bf Takeaway.}
The retained CUDA record shows host-memory mapping/copy observations after
one request and after two further requests in a partial-offload
configuration. It does not establish direct device-memory recovery,
physical-page erasure, or a universal GPU result. Fully offloaded models
and other GPU-serving architectures require separate evaluation.
%%%%%%%%%%%%%%%%%%%%%%%%%%%%%%%%%%%%%%%%%%%%%%%%%%%%%%%%%%%%%%%%%%%%%%%%%%%%%%


%%%%%%%%%%%%%%%%%%%%%%%%%%%%%%%%%%%%%%%%%%%%%%%%%%%%%%%%%%%%%%%%%%%%%%%%%%%%%%
\subsection{RQ3. Persistence Boundary}
\label{sec:rq3}
%%%%%%%%%%%%%%%%%%%%%%%%%%%%%%%%%%%%%%%%%%%%%%%%%%%%%%%%%%%%%%%%%%%%%%%%%%%%%%

RQ2 showed that prompt plaintext can outlive an inference request within
process memory. We next examine a distinct confidentiality boundary:
whether consumer wrapper applications independently extend prompt
lifetime by writing prompt-derived data to persistent storage. This distinction matters because runtime-memory sanitisation cannot
protect prompt data that have already been copied to wrapper-managed
files. We therefore ask: \emph{Do consumer local-LLM wrappers persist prompt plaintext beyond the
runtime lifetime of an inference request, and can such persistence be
disabled?}


%%%%%%%%%%%%%%%%%%%%%%%%%%%%%%%%%%%%%%%%%%%%%%%%%%%%%%%%%%%%%%%%%%%%%%%%%%%%%%

{\bf Evidence.}
We evaluate LM~Studio and Ollama under their recorded configurations by
searching retained wrapper artifacts for exact UUID matches. The LM~Studio
mitigation report identifies llmster daemon 0.0.15-2 with bundled backend
2.17.0. With \texttt{logSensitiveData=true}, 10/10 canaries were found;
with it disabled, none were found in the examined artifacts. The original
plaintext log is not retained, so this is a report-backed configuration
comparison rather than a fresh-install default claim.

The Ollama tree contains 206 metadata records: 203 trigger files completed
successfully, two contain explicit model-memory failures, and one metadata
record is incomplete without a trigger file. No Ollama executable version
is linked to these records. Six LM~Studio hook summaries show external
connections to lmstudio.ai-related hosts, but no canary in the captured or
hooked payloads; TLS was not decrypted. These observations separate local
logging from network-payload visibility.

%%%%%%%%%%%%%%%%%%%%%%%%%%%%%%%%%%%%%%%%%%%%%%%%%%%%%%%%%%%%%%%%%%%%%%%%%%%%%%

{\bf Analysis.}
These measurements show that persistent prompt lifetime is determined
independently of the runtime-memory behaviour examined in RQ2. In the
recorded LM~Studio daemon/backend configuration, enabling sensitive-data
logging produced 10/10 canary recoveries and disabling it produced none in
the examined artifacts. The original log is not retained, so the result is
configuration-specific rather than a clean-install default claim.

The Ollama accounting is descriptive: 203 trigger files completed, two
failed because of model-memory errors, and one metadata record lacks a
trigger file. The six LM~Studio hook summaries observed external
connections but no canary in visible captured payloads; because TLS was not
decrypted, this does not establish absence of all external transmission.

%%%%%%%%%%%%%%%%%%%%%%%%%%%%%%%%%%%%%%%%%%%%%%%%%%%%%%%%%%%%%%%%%%%%%%%%%%%%%%

{\bf Configuration Attribution.}
The LM~Studio intervention provides a direct configuration-level
comparison. With sensitive-data logging enabled, injected canaries remain
recoverable from wrapper-managed persistent storage; after disabling the
corresponding setting and restarting the application, the same
acquisition procedure yields no canary matches in the examined
persistent artifacts. This result bounds our conclusion to the wrapper versions,
configurations, and artifacts evaluated in this study. It does not imply
that LM~Studio can never create other prompt-bearing artifacts, nor that
Ollama never persists prompt information under every possible
configuration. Rather, it demonstrates that wrapper-level configuration
can independently determine whether prompt plaintext survives on disk
after inference.

%%%%%%%%%%%%%%%%%%%%%%%%%%%%%%%%%%%%%%%%%%%%%%%%%%%%%%%%%%%%%%%%%%%%%%%%%%%%%%

{\bf Takeaway.}
Prompt persistence is not determined solely by the inference runtime.
Under the recorded LM~Studio configuration, the logging toggle changed
whether canaries were found in the examined artifacts. The Ollama tree
provides a separate run-accounting and artifact-visibility result, not a
version-independent guarantee. External connections without visible
canaries are likewise bounded by the capture instrumentation.

%%%%%%%%%%%%%%%%%%%%%%%%%%%%%%%%%%%%%%%%%%%%%%%%%%%%%%%%%%%%%%%%%%%%%%%%%%%%%%
\subsection{RQ4. Isolation Boundary}
\label{sec:rq4}
%%%%%%%%%%%%%%%%%%%%%%%%%%%%%%%%%%%%%%%%%%%%%%%%%%%%%%%%%%%%%%%%%%%%%%%%%%%%%%

RQ2 and RQ3 examine prompt lifetime within one serving environment. We
now consider shared serving state. The retained slot experiment uses one
unauthenticated loopback server and one Python process; it does not test
two clients, distinct credentials, or ownership controls. We therefore ask:
\emph{Can saved state and subsequent continuation be observed across slot
identifiers in the evaluated shared process?} Timing evidence is treated
only as a historical, non-recomputed limitation.

%%%%%%%%%%%%%%%%%%%%%%%%%%%%%%%%%%%%%%%%%%%%%%%%%%%%%%%%%%%%%%%%%%%%%%%%%%%%%%

{\bf Same-Process Saved-State Restoration.}
When slot saving (\texttt{--slot-save-path}) is enabled, the
\texttt{/slots/:id\_slot} interface exposes operations for saving and
restoring conversation state. The retained script starts one loopback
server, sends a UUID-bearing prompt to slot~0, saves an attacker-selected
filename, and restores that filename into slot~1. It scores success only
when the reported \texttt{n\_saved} equals \texttt{n\_restored}. Across
200 trials (50 per model), this equality holds in 200/200 records. The
script has no \texttt{--api-key}, no Authorization header, no second client,
and no wrong-key or ownership control. The result is therefore same-process
unauthenticated state transfer under the tested configuration, not a
demonstrated cross-principal authorization bypass.

%%%%%%%%%%%%%%%%%%%%%%%%%%%%%%%%%%%%%%%%%%%%%%%%%%%%%%%%%%%%%%%%%%%%%%%%%%%%%%

{\bf Model-Mediated Disclosure.}
We separately test whether restored state produces a model-visible
continuation. The retained continuation files are dated July 3--4, 2026
and are not attributed here to the June campaign. They contain 126/200
KV-channel hits: Gemma 50/50, Qwen 50/50, Nemotron 25/50, and Phi-4 1/50.
The same files report 200/200 memory-channel hits. This continuation
workload is separate from the direct LM~Studio injection workload, whose
any-channel union is 150/200. The logical continuation protocol is shared,
but its retained GPU offload settings differ by model (999 offloaded layers
for Gemma, Qwen, and Nemotron versus 20 for Phi-4), so the results are not a
hardware-equivalent factorial comparison. Neither workload tests distinct
client credentials or resource ownership.

%%%%%%%%%%%%%%%%%%%%%%%%%%%%%%%%%%%%%%%%%%%%%%%%%%%%%%%%%%%%%%%%%%%%%%%%%%%%%%

{\bf Historical timing record.}
An older report describes a cached-versus-uncached timing workload, but the
per-request vectors and original analysis script are not retained. We do
not use its historical separation value as a validated quantitative result
or as evidence of arbitrary-network distinguishability.

%%%%%%%%%%%%%%%%%%%%%%%%%%%%%%%%%%%%%%%%%%%%%%%%%%%%%%%%%%%%%%%%%%%%%%%%%%%%%%

{\bf Analysis.}
The two retained workloads expose different aspects of shared-state
handling. The slot script demonstrates same-process state transfer between
slot identifiers, while the later continuation files show model- and
channel-specific recovery after restoration. Because no two-client or
two-key control was run, the records do not identify a cross-principal
authorization failure.

Model continuation is an additional model-dependent observation after
state restoration; it is not a direct test of authorization. The historical
timing record cannot be independently recomputed and is therefore omitted
from the quantitative isolation conclusion.

%%%%%%%%%%%%%%%%%%%%%%%%%%%%%%%%%%%%%%%%%%%%%%%%%%%%%%%%%%%%%%%%%%%%%%%%%%%%%%

{\bf Isolation Attribution.}
The slot equality result localises an observed state-transfer primitive to
the serving interface under one process and one credential-free protocol.
The later continuation result is a separate July-dated artifact family.
Neither provides the controls needed to infer behavior between distinct
principals.

%%%%%%%%%%%%%%%%%%%%%%%%%%%%%%%%%%%%%%%%%%%%%%%%%%%%%%%%%%%%%%%%%%%%%%%%%%%%%%

{\bf Takeaway.}
Under the evaluated one-process configuration, slot-file state transfer is
observed in 200/200 token-count equality checks. A separate July-dated
continuation workload records 126/200 KV-channel hits and 200/200
memory-channel hits. These are evidence of shared-state behavior under the
tested protocol, not demonstrations of cross-principal authorization
failure. Distinct credentials, ownership, and wrong-key controls remain
future experiments.
%%%%%%%%%%%%%%%%%%%%%%%%%%%%%%%%%%%%%%%%%%%%%%%%%%%%%%%%%%%%%%%%%%%%%%%%%%%%%%

\vspace{-0.2cm}
\section{Related Work}
\label{sec:related}
\vspace{-0.2cm}
Prior work covers local-serving vulnerabilities and multi-tenant isolation,
including malicious model parsing, template injection, authorization flaws,
and shared-architecture risks~\cite{cve_vocab_overflow,cve_bleedingllama,
cve_jinja2_ssti,cve_vllm_62164,cve_databricks,
cve_int_overflow_bypass,cve_ggml_nbytes,cve_2024_37032,cve_ollama_updater,
multitenant_llm2026,llmserving_survey2024}. Murtuza et al.~\cite{murtuza2026}
and Chung and Badhe~\cite{localnotprivate2026} document local prompt and
artifact exposure, while Xu et al.~\cite{xu2026mobilellmforensics} study
mobile-LLM forensics. Wu et al.~\cite{wu2025iknowwhatyouasked} and related
NDSS/S\&P work study shared-cache timing and reconstruction~\cite{safekv2025,
prefixwall2026,shadowcache2026,optileak2026}. Secure-memory work motivates
zeroisation and allocator controls~\cite{explicit_bzero,securestring2025,
zeroonfree2024,procfs_forensics,glibc_arena}. LLAnalyzer complements these
efforts by connecting model admission, runtime lifetime, wrapper persistence,
and shared serving state in one evidence-bounded framework.
\vspace{-0.15cm}
 %==============================================================================
 \section{Conclusion}\label{sec:conclusion}
 \vspace{-0.15cm}
 %==============================================================================
We introduced \textbf{LLAnalyzer}, a boundary-oriented framework for evaluating
prompt confidentiality across local LLM serving systems. Across the four
boundaries---\emph{Integrity}, \emph{Lifetime}, \emph{Persistence}, and
\emph{Isolation}---retained GGUF reports contain assertion-abort observations
and no recorded SIGSEGV, while runtime scans recover prompt canaries after
inference. The wrapper results are configuration-specific: LM~Studio
sensitive-data logging produced 10/10 canary recoveries in the recorded
mitigation test and disabling it produced none in the examined artifacts.
The slot workload produced 200/200 same-process token-count equality checks;
a separate July-dated continuation workload produced 126/200 KV-channel hits
and 200/200 memory-channel hits. These are not two-principal authorization
tests. Thus, \emph{locality is a deployment property, not a confidentiality
guarantee}: protecting prompts requires lifecycle-wide controls over
admission, copying, retention, and sharing, together with explicit evidence
for any claimed principal boundary.

%==============================================================================
% REFERENCES (inline; IEEE TIFS format. Every entry has a verifiable DOI / arXiv / CVE / official URL.)
%==============================================================================



% (single \appendix[...] with title kept after the bibliography; no bare \appendix here)
\section*{Open Science}
\label{app:open_science}
%==============================================================================

The project reports record source revision 721354f (full SHA
721354fbdfb7743e2be2183d918a3cdb9276c70f). The binaries executed in June
2026 were not hashed and are no longer available. An October 2026 CPU-only
rebuild is retained for reference, is not claimed byte-identical, and its
equivalence to the executed binary is unknown. LM~Studio runs used its
bundled llama.cpp backend; the GPU runs used a CUDA build whose revision is
not recorded. A public artifact URL remains \artifacturl{TODO}; the
retained records include scripts, reports, and selected raw evidence, but
not every historical binary or timing vector.


%==============================================================================
\section*{Ethical Considerations}
\label{app:ethics}
%==============================================================================

All experiments used researcher-controlled systems, accounts, tenants, and
synthetic UUID-tagged prompts; no human-subject data, third-party conversations,
production systems, or public-host state were accessed.

Internet measurements were passive and non-destructive. Observations involving
shared-serving behaviour are reported as operator-managed deployment-hardening
guidance. No vendor report has been transmitted and no vendor decision is
established; the dated scope statement is researcher-side.
\bibliographystyle{IEEEtran}
\bibliography{references}

\end{document}
